% !TeX root = ../../paper.tex

Modern VR headsets are particularly effective at stimulating the senses of sight and hearing but they offer no resistance when the player reaches out. This limitation becomes especially apparent in the moments when a VR experience works hardest to create an immersive situation. The player leans in, stretches out a hand, and touches nothing. Force-feedback hardware and haptic suits could close the gap, but they remain expensive, cumbersome, and restrict free movement. Passive haptics skips the hardware. It involves placing real objects at the locations where the virtual objects appear to be, so that the player's hand finds what their eyes expect~\cite{Insko.2001, substitutional-reality}.

The horror genre is built on anticipation: the player knows something unpleasant is coming and has to act anyway. Reaching into a bowl one cannot see into is unsettling on a screen; when the hand sinks into something cold and yielding, the stakes change. Horror also punishes broken illusions, because a scare loses its force the moment a hand swipes through empty air. Presence and fear are coupled in VR: in virtual height exposure the two correlate at $r = .62$~\cite{Gromer.2019}, and adding a physical ledge to a virtual pit raised participants' heart-rate response as well as their reported presence~\cite{Meehan.2002}. If passive haptics changes presence at all, an experience built on disgust and fear should make the difference visible.

\emph{Tangible Horror}, our VR escape scenario for the Meta Quest~3, puts this to the test. The game locks the player in a run-down cabin; to open the lock on the door, they must collect four coins, each hidden inside a small horror challenge that the player has to reach into.
In the haptic condition each challenge is supplemented by a low-cost physical proxy (rubber balls, gummy worms, a water-spraying tube operated by an experimenter, and synthetic Halloween webbing) aligned with its virtual counterpart. In the baseline condition the player plays the identical scene and touches nothing but air. We compared both conditions within participants, in randomised order, using the Igroup Presence Questionnaire (IPQ)~\cite{Schubert.2001} after each playthrough.

This paper makes two contributions. First, we describe the design and implementation of the four proxy challenges, including the practical work that decides whether a proxy illusion holds: calibrating virtual object sizes to their physical stand-ins, anchoring furniture in the tracking space, and tracking hand-held coins with fiducial markers. Second, we report an empirical comparison of the haptic and non-haptic versions of the same experience. Section~\ref{sec:related-work} situates the project in prior work on passive haptics and on presence and emotion in VR. Sections~\ref{sec:apparatus} and~\ref{sec:study} describe the prototype and the study; Sections~\ref{sec:results} and~\ref{sec:discussion} present and discuss the results.
